\section{A vanishing theorem for sparse simplicial cycles}
\label{sec:framework}

We formulate the geometric statement independently of group rings.
Its input is a graph with a free group action and a simplicial complex
whose top-dimensional vertex sets are connected in that graph. Its output
is the vanishing of every invariant cocycle evaluation in sufficiently
high dimension.

For a set $Y$, let $\C^{(Y)}$ be the vector space with basis
$\{\delta_y:y\in Y\}$, consisting of finite linear combinations. The
augmented chains of the full simplex on $Y$ are
\[
 E_j=\bigwedge^{j+1}\C^{(Y)}\quad(j\geq0),\qquad E_{-1}=\C.
\]
The differential is contraction with the functional
$\alpha(\sum_y c_y\delta_y)=\sum_y c_y$:
\begin{equation}\label{eq:boundary}
 \partial(\delta_{y_0}\wedge\cdots\wedge\delta_{y_j})
 =\sum_{i=0}^{j}(-1)^i
 \delta_{y_0}\wedge\cdots\wedge\widehat{\delta_{y_i}}
 \wedge\cdots\wedge\delta_{y_j}.
\end{equation}
In degree zero this is the augmentation to $\C$; the empty wedge is $1$.
For a simplicial subcomplex $K$, its augmented oriented chains
$C_*(K)$ form a subcomplex of $E_*$.

If a group $H$ acts on $Y$, it acts on both chain complexes by permuting
vertices, with the signs prescribed by the exterior product. For an
$H$-module $V$, its coinvariants are
\[
 V_H=V/\operatorname{span}_{\C}\{hv-v:h\in H,\ v\in V\}.
\]
The differentials induce differentials on coinvariants. A cycle in
$C_n(K)_H$ is a class whose boundary is zero there; an arbitrary finite
chain representing it need not have zero boundary before passage to
coinvariants. An $H$-invariant linear functional $\psi:E_n\to\C$
descends to $(E_n)_H$. We call it a cocycle when
$\psi\partial=0$ on $E_{n+1}$.

\begin{theorem}[Sparse-chain vanishing]\label{thm:sparse}
For each integer $D\geq2$ there is an integer $n_0(D)$ with the following
property. Let $Y$ be the vertex set of a countable connected simple graph
of degree at most $D$, and let a countable group $H$ act by graph
automorphisms, freely on vertices. Suppose $n\geq n_0(D)$ and $K$ is an
$H$-invariant subcomplex of the full simplex on $Y$ such that:
\begin{enumerate}[label=\textup{(\roman*)}]
 \item $K$ is generated by its $n$-simplices and has finitely many
 $H$-orbits of simplices;
 \item the vertices of each $n$-simplex span a connected subgraph of $Y$.
\end{enumerate}
Then every $H$-invariant cocycle $\psi:E_n\to\C$ evaluates to zero on
the image in $(E_n)_H$ of every cycle in $C_n(K)_H$.
\end{theorem}

The bound $n_0(D)$ depends only on the graph degree, not on $H$, $K$ or
$\psi$. There is no boundedness hypothesis on $\psi$. We will bound it
only on a finite set of basis wedges associated with each fixed original
simplex. This distinction is necessary in the application.

For the proof assume $K\ne\varnothing$; the empty case has no positive
degree chains. Its realization $|K|$ consists of finitely supported
probability vectors $p:Y\to[0,1]$ whose supports are simplices of $K$.
A simplex will also denote its closed affine realization; $\sigma^\circ$
denotes its relative interior. We call simplices of $K$ \emph{original}
to distinguish them from the later subdivisions, and call its
$n$-simplices \emph{top simplices}.

The proof has three steps. First, we choose a metric and simplexwise
areas for which balls have small volume in high dimension.
Second, we construct separating filtrations with very few points in
their final stratum, on average over a free probability parameter.
Third, we use their subdivisions to construct equivariant chain maps
homotopic to inclusion whose cocycle evaluations approach zero. Each
approximation uses a finite filtration; no limiting filtration is needed.
